%%==================================================
%% abstract.tex for XMU Graduate Thesis
%%==================================================

\begin{abstract}

碳纳米管阵列因优良的电学、热学和力学性能，在储能、热管理和高性能复合材料等领域具有重要应用潜力。围绕阵列制备与应用开发，工艺条件如何影响形貌演化并进一步作用于结构状态和性能表现，是相关研究中的关键问题。现有研究主要面临两方面困难。一方面，扫描电子显微镜（SEM）图像是分析阵列整体形貌的主要依据，但碳纳米管阵列在图像中表现为细长交叠结构，样品级形貌特征难以稳定量化，工艺—形貌分析因而缺少可靠的数据支撑。另一方面，碳纳米管阵列工艺—形貌关系具有较强的条件依赖性和机理复杂性，仅依赖样品统计数据难以支撑关键变量作用解释与结果追溯，而工艺条件、形貌特征和生长机理等信息又长期分散于文献、实验记录和图像分析结果中，难以形成有效的知识支撑。受此影响，碳纳米管阵列工艺—形貌关系研究仍缺少稳定且可追溯的分析基础。

本文首先建立多源知识组织与关系约束检索方法。针对学术文献、实验记录、工艺参数表、结构化实验数据及图像分析结果长期分散且难以直接调用的问题，建立统一的知识组织与关系约束检索方法，将相关信息整理为可关联的知识单元，为工艺变量筛选和机理分析提供证据基础。

随后建立基于 CNTSegNet-clDice 的样品级自动化形貌表征方法。针对碳纳米管阵列 SEM 图像中复杂细长结构前景提取困难的问题，构建兼顾连通性保持的前景分割方法，并结合骨架分析与几何计算建立样品级量化流程，提取密度、取向度、表观直径、平均曲率、波曲度和迂曲度等形貌特征，通过与人工分级结果比较验证样品级形貌标签的有效性。

在此基础上，本文结合知识检索结果与样品级形貌指标，对实验制备和表征样品中的关键工艺变量开展工艺—形貌关系分析，并选取 Fe 溅射功率、Fe 催化层厚度和退火时间进行统计分析和随机森林建模。结果表明，在 Fe 磁控溅射催化层和热 CVD 生长构成的样品体系中，三类变量与取向度、平均曲率、波曲度和迂曲度之间存在可建模的对应关系。其中，退火时间与取向度下降及弯曲相关特征增强之间呈现最稳定的趋势关系，Fe 溅射功率主要影响不同工艺区间中的响应强弱及其出现区间，Fe 催化层厚度则表现为相对较弱的补充调节作用。结合知识证据可知，三类变量主要通过催化层形成与颗粒演化共同影响阵列取向和弯曲响应。相关方法已在科研辅助分析平台中完成集成验证，表明该分析流程能够支持碳纳米管阵列样品级工艺—形貌关系的定量分析与机理解释。

\end{abstract}

\keywords{碳纳米管阵列；碳纳米管；知识增强；形貌表征；工艺—形貌关系}

\begin{enabstract}

Carbon nanotube (CNT) arrays have broad application prospects in energy storage devices, thermal interface materials, field emission devices, and high-performance composites because of their excellent electrical, thermal, and mechanical properties. A central issue in their preparation and application is how process conditions drive morphology evolution and further affect structural state and performance. Existing studies mainly face two difficulties. First, scanning electron microscopy (SEM) images are the primary basis for characterizing overall array morphology, yet the slender and overlapping structures in these images make stable quantification of sample-level morphology features difficult, leaving process-morphology analysis without reliable data support. Second, process-morphology relationships in CNT arrays are strongly condition-dependent and mechanistically complex. Sample statistics alone are insufficient for explaining the effects of key variables and supporting traceable analysis, while information on process conditions, morphology features, and growth mechanisms remains scattered across literature, experimental records, and image analysis results. As a result, studies of CNT array process-morphology relationships still lack a stable and traceable analytical basis.

To address this problem, the thesis first develops a multi-source knowledge organization and relation-constrained retrieval method. For academic literature, experimental records, process-parameter tables, structured experimental data, and image analysis results that are scattered and difficult to use directly, a unified organization scheme is established to turn relevant information into linkable knowledge units and provide evidence for process-variable screening and mechanism analysis.

A sample-level automated morphology characterization method, CNTSegNet-clDice, is then developed. To address the foreground extraction problem posed by complex slender structures in SEM images of CNT arrays, a segmentation method that preserves structural connectivity is proposed. Combined with skeleton analysis and geometric computation, it forms a sample-level quantification pipeline for extracting density, alignment, apparent diameter, average curvature, waviness, and tortuosity. The resulting sample-level morphology labels are validated through comparison with manual grading results.

On this basis, the thesis investigates process-morphology relationships by combining knowledge retrieval results with sample-level morphology indicators, and conducts statistical analysis and random forest modeling for experimentally prepared and characterized samples, with Fe sputtering power, Fe catalyst layer thickness, and annealing time selected as the key process variables. The results show that, in the sample system defined by Fe magnetron-sputtered catalyst layers and thermal CVD growth, these three variables have modelable correspondences with alignment, average curvature, waviness, and tortuosity. Among them, annealing time shows the most stable trend relationship with decreasing alignment and enhanced bending-related features, Fe sputtering power mainly affects response intensity and the corresponding annealing ranges, and Fe catalyst layer thickness plays a relatively weak supplementary regulatory role. Knowledge evidence further indicates that these three variables jointly affect array alignment and bending response through catalyst-layer formation and particle evolution. The method has been integrated and validated in a scientific research assistant platform, indicating that the workflow supports quantitative sample-level process-morphology analysis and mechanistic interpretation for CNT arrays.

\end{enabstract}

\enkeywords{Carbon Nanotube Arrays; Carbon Nanotubes; Knowledge Enhancement; Morphology Characterization; Process-Morphology Relationship}
